%==============================================================
\section{Cơ sở điện toán lượng tử}
\label{sec:basics}
%==============================================================

Chương này trả lời ba câu hỏi nền tảng: một qubit được mô tả bằng gì, \emph{superposition} thực chất
là gì, và phép đo làm gì với trạng thái đó. Phần cuối mở rộng sang hệ nhiều qubit.

\subsection{Qubit và ký hiệu Dirac}

Bit cổ điển chỉ nhận một trong hai giá trị $0$ hoặc $1$. Qubit là một \textbf{vector đơn vị} trong
không gian Hilbert phức hai chiều $\mathcal{H} = \mathbb{C}^2$. Hai vector cơ sở được ký hiệu
\begin{equation}
  \ket{0} = \begin{pmatrix}1\\0\end{pmatrix}, \qquad
  \ket{1} = \begin{pmatrix}0\\1\end{pmatrix}.
  \label{eq:basis}
\end{equation}

\begin{definition}[Ket, bra và tích trong]
Ký hiệu $\ket{\psi}$ (đọc là ``ket psi'') là một \textbf{vector cột} biểu diễn trạng thái lượng tử.
Ký hiệu $\bra{\psi} = \ket{\psi}^\dagger$ (``bra psi'') là \textbf{vector hàng} liên hợp phức của
nó. Ghép bra với ket cho \textbf{tích trong} (inner product) $\braket{\phi|\psi} \in \mathbb{C}$;
ghép ket với bra cho \textbf{tích ngoài} (outer product) $\ket{\psi}\bra{\phi}$, là một ma trận.
Cách viết này gọi là \textbf{ký hiệu Dirac} (bra--ket notation) \cite{nielsen2010}.
\end{definition}

Hai vector $\ket{0}, \ket{1}$ trực giao ($\braket{0|1} = 0$) và có chuẩn bằng 1
($\braket{0|0} = \braket{1|1} = 1$); chúng tạo thành \textbf{computational basis} (cơ sở tính
toán, còn gọi là cơ sở $Z$). Trong Qiskit, mọi qubit luôn được khởi tạo ở $\ket{0}$.

\subsection{Superposition và quy tắc Born}
\label{sec:born}

Trạng thái tổng quát của một qubit là tổ hợp tuyến tính có chuẩn hoá của hai trạng thái cơ sở:
\begin{equation}
  \ket{\psi} = \alpha\ket{0} + \beta\ket{1}, \qquad
  \alpha, \beta \in \mathbb{C}, \qquad |\alpha|^2 + |\beta|^2 = 1.
  \label{eq:superposition}
\end{equation}
Các hệ số $\alpha, \beta$ gọi là \textbf{amplitude} (biên độ xác suất). Khi đo trong cơ sở $Z$, theo
\textbf{quy tắc Born}, phép đo trả về \textbf{bit cổ điển} $0$ với xác suất
$P(0) = |\alpha|^2 = |\braket{0|\psi}|^2$, và khi đó qubit chuyển về trạng thái $\ket{0}$; hoặc trả
về bit $1$ với xác suất $P(1) = |\beta|^2$, và qubit chuyển về $\ket{1}$. Cần phân biệt
\textbf{kết quả đo} (một bit cổ điển, ví dụ chuỗi \code{'0'} trong counts của Qiskit) với
\textbf{trạng thái sau khi đo} (vector $\ket{0}$). Điều kiện chuẩn hoá bảo đảm tổng xác suất bằng 1.

Ba trạng thái superposition thường gặp:
\begin{equation}
  \ket{+} = \tfrac{1}{\sqrt2}\big(\ket{0} + \ket{1}\big), \qquad
  \ket{-} = \tfrac{1}{\sqrt2}\big(\ket{0} - \ket{1}\big), \qquad
  \ket{i} = \tfrac{1}{\sqrt2}\big(\ket{0} + i\ket{1}\big).
  \label{eq:plusminus}
\end{equation}
Cả ba đều cho $P(0) = P(1) = 0.5$ khi đo trong cơ sở $Z$, nhưng chúng \textbf{không phải} cùng một
trạng thái -- sự khác biệt nằm ở \emph{phase} của amplitude (Mục~\ref{sec:phase}).

\begin{keybox}
Superposition \textbf{không} có nghĩa là ``qubit vừa là 0 vừa là 1'' theo nghĩa cổ điển. Nó có nghĩa
là trạng thái là một vector có hai thành phần phức, và hai thành phần này có thể \emph{giao thoa}
(interference) với nhau dưới tác dụng của các cổng: cộng dồn hoặc triệt tiêu lẫn nhau. Chính sự giao
thoa này mới là nguồn gốc lợi thế của thuật toán lượng tử.
\end{keybox}

\subsection{Bloch sphere}

Bỏ qua global phase (xem Mục~\ref{sec:phase}), mọi trạng thái một qubit đều viết được dưới dạng
\begin{equation}
  \ket{\psi} = \cos\tfrac{\theta}{2}\ket{0} + e^{i\varphi}\sin\tfrac{\theta}{2}\ket{1},
  \qquad 0 \le \theta \le \pi,\ 0 \le \varphi < 2\pi.
  \label{eq:bloch}
\end{equation}
Cặp góc $(\theta, \varphi)$ xác định một điểm trên mặt cầu đơn vị gọi là \textbf{Bloch sphere}
(Hình~\ref{fig:bloch}). Các trạng thái $\ket{0}, \ket{1}$ nằm ở hai cực, $\ket{\pm}$ trên trục $X$,
$\ket{\pm i}$ trên trục $Y$. Đáng chú ý: hai trạng thái \emph{trực giao} như $\ket{0}$ và $\ket{1}$
nằm \textbf{đối diện} nhau trên cầu chứ không vuông góc. Trên Bloch sphere, mọi cổng một qubit là một
\textbf{phép quay}.

\begin{figure}[h]
  \centering
  \begin{tikzpicture}[scale=2.1, >=Stealth]
    % sphere
    \shade[ball color=AccentBlue!8, opacity=0.6] (0,0) circle (1);
    \draw[thick] (0,0) circle (1);
    \draw[dashed, black!50] (1,0) arc (0:180:1 and 0.3);
    \draw[black!70] (-1,0) arc (180:360:1 and 0.3);
    % axes
    \draw[->] (0,0) -- (0,1.25) node[above] {$z$};
    \draw[->] (0,0) -- (-0.55,-0.55) node[below left] {$x$};
    \draw[->] (0,0) -- (1.3,0) node[right] {$y$};
    \draw[dashed, black!40] (0,0) -- (0,-1);
    % labels of basis states
    \fill (0,1) circle (0.025) node[above right] {$\ket{0}$};
    \fill (0,-1) circle (0.025) node[below right] {$\ket{1}$};
    \fill (-0.42,-0.42) circle (0.025) node[left=2pt] {$\ket{+}$};
    \draw[dashed, black!35] (0,0) -- (0.42,0.42);
    \fill[black!60] (0.42,0.42) circle (0.02) node[right=1pt, black!60] {$\ket{-}$};
    \fill (1,0) circle (0.025) node[below right] {$\ket{i}$};
    % state vector (front-right, away from the -x direction)
    \coordinate (P) at (0.35,0.65);
    \coordinate (Pp) at (0.35,-0.25);
    \draw[->, very thick, ProbRed] (0,0) -- (P) node[above left=-1pt] {$\ket{\psi}$};
    \draw[dashed, ProbRed!70] (P) -- (Pp);
    \draw[dashed, ProbRed!70] (0,0) -- (Pp);
    % angles
    \draw[OkGreen, thick] (0,0.32) arc (90:62:0.32);
    \node[OkGreen] at (0.08,0.45) {$\theta$};
    \draw[Purple, thick] (-0.2,-0.2) arc (225:325:0.28 and 0.12);
    \node[Purple] at (0.05,-0.38) {$\varphi$};
  \end{tikzpicture}
  \caption{Bloch sphere. Trạng thái $\ket{\psi}$ ở phương trình~\eqref{eq:bloch} ứng với điểm có
  góc cực $\theta$ và góc phương vị $\varphi$.}
  \label{fig:bloch}
\end{figure}

\subsection{Phép đo và collapse}

Như đã nêu ở Mục~\ref{sec:born}, sau phép đo trạng thái \textbf{collapse} (sụp đổ) về đúng ket
tương ứng với bit đo được: kết quả $0$ ứng với $\ket{0}$, kết quả $1$ ứng với $\ket{1}$. Hệ quả thực
tế:
\begin{itemize}[noitemsep]
  \item Một lần đo \textbf{không} cho biết $\alpha, \beta$. Muốn ước lượng phân bố xác suất phải chuẩn
    bị lại trạng thái và đo nhiều lần; mỗi lần gọi là một \textbf{shot}.
  \item Trong Qiskit 2.x, primitive \code{SamplerV2} trả về số lần xuất hiện (\emph{counts}) của mỗi
    chuỗi bit; \code{Statevector} trả về amplitude chính xác nhưng chỉ dùng được với mạch
    \emph{không có} lệnh đo.
\end{itemize}

\subsection{Relative phase và global phase}
\label{sec:phase}

\paragraph{Relative phase.} Hai trạng thái $\ket{+}$ và $\ket{-}$ chỉ khác nhau ở \textbf{dấu của hệ
số đứng trước $\ket{1}$} ($+\tfrac{1}{\sqrt2}$ so với $-\tfrac{1}{\sqrt2}$). Tổng quát, đại lượng phân
biệt chúng là độ lệch pha $\varphi$ giữa $\beta$ và $\alpha$, gọi là \textbf{relative phase} (pha
tương đối); dấu trừ ứng với $\varphi = \pi$, vì $e^{i\pi} = -1$. Tương tự, $\ket{i}$ có
$\varphi = \pi/2$ vì $e^{i\pi/2} = i$. Góc $\varphi$ này chính là góc phương vị trên Bloch sphere ở
phương trình~\eqref{eq:bloch}. Đo trong cơ sở $Z$ không phân biệt được $\ket{+}$ và $\ket{-}$, vì
$|\beta|^2 = \tfrac12$ trong cả hai trường hợp; nhưng đổi cơ sở đo thì được: áp cổng Hadamard trước
khi đo, vì $H\ket{+} = \ket{0}$ và $H\ket{-} = \ket{1}$. Relative phase vì vậy là đại lượng
\textbf{quan sát được}.

\paragraph{Global phase.} \textbf{Global phase} (pha toàn cục) là trường hợp nhân \emph{cả} vector
trạng thái -- tức \emph{mọi} amplitude cùng lúc -- với cùng một hệ số $e^{i\gamma}$:
\begin{equation}
  \ket{\psi'} = e^{i\gamma}\ket{\psi} = e^{i\gamma}\alpha\ket{0} + e^{i\gamma}\beta\ket{1}.
  \label{eq:global}
\end{equation}
Vì mọi amplitude bị quay pha như nhau, độ lệch pha giữa chúng -- relative phase -- không đổi. Hai
lập luận cho thấy global phase \textbf{không thể quan sát được} bằng bất kỳ thí nghiệm nào:
\begin{enumerate}[noitemsep]
  \item \textbf{Mọi phép đo cho cùng phân bố xác suất.} Với phép đo trong một cơ sở trực chuẩn bất kỳ
    $\{\ket{\phi_k}\}$ (không chỉ cơ sở $Z$), quy tắc Born cho
    \[
      P'(k) = \big|\braket{\phi_k|\psi'}\big|^2 = |e^{i\gamma}|^2\,\big|\braket{\phi_k|\psi}\big|^2
      = \big|\braket{\phi_k|\psi}\big|^2 = P(k),
    \]
    vì $|e^{i\gamma}| = 1$.
  \item \textbf{Không cổng nào biến nó thành relative phase.} Với mọi unitary $U$,
    $U\big(e^{i\gamma}\ket{\psi}\big) = e^{i\gamma}\,U\ket{\psi}$: hệ số $e^{i\gamma}$ đi xuyên qua cả
    mạch mà vẫn là hệ số chung, nên đến lúc đo nó vẫn không ảnh hưởng gì (theo lập luận 1).
\end{enumerate}
Vì vậy $\ket{\psi}$ và $e^{i\gamma}\ket{\psi}$ là \textbf{cùng một trạng thái vật lý}. Ví dụ,
$-\ket{-} = \tfrac{1}{\sqrt2}(-\ket{0} + \ket{1})$ chính là $\ket{-}$: cả hai hệ số đều đổi dấu, nên
hệ số trước $\ket{1}$ vẫn ngược dấu với hệ số trước $\ket{0}$. Ngược lại, $\ket{+}$ và $\ket{-}$ là
hai trạng thái khác nhau, vì chỉ \emph{một} hệ số đổi dấu.

Global phase có ba hệ quả được dùng xuyên suốt báo cáo:
\begin{itemize}[noitemsep]
  \item \textbf{Bloch sphere chỉ cần hai góc.} Một qubit có hai amplitude phức (4 tham số thực); trừ
    điều kiện chuẩn hoá và global phase, còn lại đúng 2 tham số. Đây là lý do phương
    trình~\eqref{eq:bloch} được phép chọn hệ số của $\ket{0}$ là số thực không âm $\cos\tfrac{\theta}{2}$.
  \item \textbf{Hai cổng sai khác global phase là cùng một cổng.} Ví dụ $R_X(\pi) = -iX$ và
    $R_Z(\lambda) = e^{-i\lambda/2}P(\lambda)$ (Mục~\ref{sec:gates}). Kết quả mô phỏng có thể khác
    công thức một dấu chung -- như trạng thái Bell $\ket{\Psi^-}$ xuất hiện dưới dạng
    $\tfrac{1}{\sqrt2}(\ket{10}-\ket{01})$ -- mà vẫn đúng. Đây là lý do Qiskit phân biệt
    \code{Operator.equiv} (bỏ qua global phase) với phép so sánh \code{==} (so từng phần tử ma trận).
  \item \textbf{Ngoại lệ quan trọng: cổng controlled.} Global phase của một cổng $U$ không quan sát
    được khi $U$ được áp \emph{vô điều kiện}, nhưng trở thành \emph{relative} phase khi $U$ chỉ được
    áp lúc qubit control bằng 1. Chẳng hạn, áp $-I$ lên target có điều kiện thì chỉ thành phần
    $\ket{1}$ của control bị đổi dấu -- tức tương đương cổng $Z$ trên control. Đây chính là cơ chế
    phase kickback ở Mục~\ref{sec:algorithms}, và là lý do khi tổng hợp cổng controlled, compiler phải
    theo dõi global phase (thuộc tính \code{QuantumCircuit.global\_phase} trong Qiskit) thay vì bỏ nó đi.
\end{itemize}

\begin{keybox}
Relative phase là ``nguyên liệu'' của mọi thuật toán lượng tử: thuật toán ghi thông tin vào phase,
rồi dùng interference (thường là cổng $H$) để biến phase thành xác suất đo được. Global phase thì không
mang thông tin vật lý -- trừ khi nó bị ``gắn điều kiện'' qua một cổng controlled.
\end{keybox}

\subsection{Hệ nhiều qubit và tích tensor}

Hệ gồm hai qubit được mô tả trong không gian $\mathbb{C}^2 \otimes \mathbb{C}^2 = \mathbb{C}^4$.
Hai qubit độc lập ghép lại bằng \textbf{tích tensor} (tích Kronecker):
\begin{equation}
  \ket{a}\otimes\ket{b} =
  \begin{pmatrix}a_0\\a_1\end{pmatrix} \otimes \begin{pmatrix}b_0\\b_1\end{pmatrix}
  = \begin{pmatrix}a_0b_0\\a_0b_1\\a_1b_0\\a_1b_1\end{pmatrix}.
\end{equation}
Tổng quát, $n$ qubit cần $2^n$ amplitude: $\ket{\psi} = \sum_{x \in \{0,1\}^n} c_x \ket{x}$ với
$\sum_x |c_x|^2 = 1$. Không gian trạng thái tăng theo \emph{hàm mũ} -- đây vừa là nguồn sức mạnh của
máy tính lượng tử, vừa là lý do mô phỏng cổ điển nhanh chóng trở nên bất khả thi.

\subsection{Quy ước thứ tự qubit trong Qiskit (little-endian)}
\label{sec:endian}

Khi ghép $n$ qubit bằng tích tensor, phải chọn qubit nào đứng \emph{trái} và qubit nào đứng
\emph{phải} trong ket. Có hai quy ước, và Qiskit dùng quy ước \emph{ngược} với phần lớn giáo trình:
\begin{itemize}[noitemsep]
  \item \textbf{Giáo trình} (Nielsen \& Chuang \cite{nielsen2010}): viết $\ket{q_0 q_1 \dots q_{n-1}}$
    và $\ket{\psi} = \ket{\psi_0}\otimes\ket{\psi_1}\otimes\dots\otimes\ket{\psi_{n-1}}$ -- qubit 0
    đứng \textbf{bên trái}.
  \item \textbf{Qiskit}: viết $\ket{q_{n-1}\dots q_1 q_0}$ và
    $\ket{\psi} = \ket{\psi_{n-1}}\otimes\dots\otimes\ket{\psi_1}\otimes\ket{\psi_0}$ -- qubit 0 đứng
    \textbf{bên phải}.
\end{itemize}

Ở cả hai quy ước, thứ tự viết trong ket luôn trùng với thứ tự các thừa số của tích Kronecker từ trái
sang phải; chỉ có việc qubit nào đứng trái là khác nhau. Nói gọn: \textbf{trong Qiskit, qubit có chỉ
số lớn hơn đứng trước} (bên trái), cả trong ket lẫn trong tích tensor. Nếu ghép theo thứ tự
$\ket{q_0}\otimes\ket{q_1}\otimes\dots$ như giáo trình rồi đưa vào Qiskit, ta được vector bị đảo thứ tự
qubit: với $q_0 = 1, q_1 = 0$, tích $\ket{q_1}\otimes\ket{q_0} = (0,1,0,0)^{\mathsf T}$ là đúng
(chuỗi \code{'01'}), còn $\ket{q_0}\otimes\ket{q_1} = (0,0,1,0)^{\mathsf T}$ lại là trạng thái có
qubit 1 bằng 1.

\paragraph{Vì sao gọi là little-endian.} Mỗi ket cơ sở $\ket{q_{n-1}\dots q_1 q_0}$ ứng với một chỉ số
nguyên $k$ trong vector trạng thái $2^n$ phần tử. Qiskit đọc chuỗi bit đó như cách viết số nhị phân
thông thường, nên
\begin{equation}
  k = \sum_{i=0}^{n-1} q_i\, 2^i = q_0\cdot 2^0 + q_1\cdot 2^1 + \dots + q_{n-1}\cdot 2^{n-1}.
  \label{eq:endian}
\end{equation}
Qubit 0 là \textbf{bit có trọng số thấp nhất} (least significant bit) của chỉ số $k$ -- giống thứ tự
byte ``little-endian'' trong kiến trúc máy tính, nơi phần có trọng số thấp đứng trước. Hệ quả là
qubit 0 xuất hiện ở \textbf{ký tự bên phải cùng} của mọi chuỗi bit mà Qiskit in ra.

\paragraph{Ví dụ.} Mạch 3 qubit, chỉ áp $X$ lên qubit 0:
\begin{itemize}[noitemsep]
  \item Qiskit trả về trạng thái $\ket{001}$, tức chỉ số $k = 1$; vector trạng thái là
    $(0,1,0,0,0,0,0,0)^{\mathsf T}$, và counts là \code{\{'001': ...\}}.
  \item Giáo trình viết cùng trạng thái đó là $\ket{100}$.
\end{itemize}
Hai cách viết mô tả \emph{cùng một} trạng thái vật lý; chỉ có thứ tự ghi các qubit là khác nhau.
Tương tự, áp $H$ lên qubit 0 của mạch 2 qubit cho $\tfrac{1}{\sqrt2}(\ket{00}+\ket{01})$ theo Qiskit
(vector $(1,1,0,0)^{\mathsf T}/\sqrt2$), nhưng là $\tfrac{1}{\sqrt2}(\ket{00}+\ket{10})$ theo giáo trình.

\paragraph{Ảnh hưởng tới ma trận cổng.} Vì thứ tự tensor bị đảo, ma trận của cổng nhiều qubit trong
Qiskit khác ma trận trong giáo trình. Cổng $H$ trên qubit 0 có ma trận $I\otimes H$ (qubit 0 là thừa
số bên phải). Cổng CX với control là qubit 0 và target là qubit 1 (\code{qc.cx(0, 1)}) có ma trận
\[
  \mathrm{CX}^{\text{Qiskit}}_{0\to1} =
  \begin{pmatrix}1&0&0&0\\0&0&0&1\\0&0&1&0\\0&1&0&0\end{pmatrix}
  \qquad\text{thay vì}\qquad
  \mathrm{CX}^{\text{giáo trình}} =
  \begin{pmatrix}1&0&0&0\\0&1&0&0\\0&0&0&1\\0&0&1&0\end{pmatrix}.
\]
Ma trận của Qiskit đổi chỗ $\ket{01}\leftrightarrow\ket{11}$ (các trạng thái có $q_0 = 1$), còn ma
trận giáo trình đổi chỗ $\ket{10}\leftrightarrow\ket{11}$ (các trạng thái có qubit đứng trái bằng 1).
Cả hai đều đúng theo quy ước của mình.

\paragraph{Ảnh hưởng tới kết quả đo.} Chuỗi counts có dạng $c_{n-1}\dots c_1 c_0$: bit cổ điển 0 đứng
bên phải. Ví dụ, đo mạch 3 qubit đã áp $X$ lên qubit 0 và 1 (với \code{measure([0,1,2],[0,1,2])}) cho
chuỗi \code{'011'}. Lưu ý rằng khi \emph{vẽ} mạch, Qiskit lại đặt qubit 0 ở dòng \emph{trên cùng} --
nên thứ tự từ trên xuống trong hình vẽ ngược với thứ tự từ trái sang phải trong chuỗi kết quả.

\paragraph{Cách làm việc an toàn.}
\begin{itemize}[noitemsep]
  \item Khi đọc kết quả, ký tự thứ $i$ \emph{tính từ bên phải} (\code{s[-1-i]} trong Python) là giá trị
    của qubit/bit cổ điển $i$. Muốn đọc theo thứ tự giáo trình thì đảo chuỗi (\code{s[::-1]}).
  \item Khi so sánh với công thức trong giáo trình, dùng \code{QuantumCircuit.reverse\_bits()} (đảo thứ
    tự qubit của mạch) hoặc \code{Statevector.reverse\_qargs()}: mạch ở ví dụ trên sau
    \code{reverse\_bits()} cho đúng \code{'100'}.
  \item Khi ghép trạng thái bằng tích tensor trong Qiskit, phải viết \code{q1.tensor(q0)} (qubit có chỉ
    số lớn đứng trước).
\end{itemize}
Quy ước này xuất hiện ở mọi thuật toán của Mục~\ref{sec:algorithms}: ký tự cuối của chuỗi bí mật $s$
trong Bernstein--Vazirani ứng với qubit 0, và chuỗi kết quả của superdense coding bị đảo so với thứ tự
viết tin nhắn.

\begin{table}[h]
  \centering
  \caption{Tóm tắt các khái niệm cơ bản và cách biểu diễn trong Qiskit.}
  \label{tab:basics}
  \small
  \begin{tabular}{@{}lll@{}}
    \toprule
    \textbf{Khái niệm} & \textbf{Ý nghĩa} & \textbf{Qiskit} \\
    \midrule
    Ket $\ket{0}, \ket{1}$ & computational basis & \code{Statevector.from\_label("0")} \\
    Amplitude $\alpha, \beta$ & số phức, $|\alpha|^2+|\beta|^2=1$ & \code{sv.data}, \code{sv.is\_valid()} \\
    Xác suất & $P(k) = |c_k|^2$ (quy tắc Born) & \code{sv.probabilities\_dict()} \\
    Superposition & tổ hợp tuyến tính các ket & \code{qc.h(0)} \\
    Đo & lấy mẫu + collapse & \code{qc.measure(...)} + \code{SamplerV2} \\
    Relative phase & quan sát được (đổi cơ sở đo) & thêm \code{h} trước \code{measure} \\
    Global phase & không quan sát được & \code{Operator.equiv} \\
    Thứ tự bit & little-endian, qubit 0 bên phải & \code{'001'}: qubit 0 bằng 1 \\
    \bottomrule
  \end{tabular}
\end{table}
